\subsection{Breakdown as a physical prediction}
  \label{subsec:breakdown-as-a-physical-prediction}

  Reconstruction fails when (i) connectivity becomes highly non-local or
  (ii) the spectrum of $G$ exhibits no clear eigenvalue gap
  (glassy or fractal regimes).
  This signals a transition to a pre-geometric regime in which
  a smooth manifold description ceases to be valid.

  \paragraph{Status of the effective metric.}
    The metric arises operationally from relational propagation and
    coarse-grained relaxation properties.
    It encodes how relational distinctions map onto effective spatial
    and temporal separations, but the mapping is many-to-one:
    distinct relational configurations may yield identical metrics,
    and relational changes need not modify the effective geometry.

  \paragraph{Regime-dependent field equations.}
    Poisson-type and wave-like equations emerge from linearizing the
    projected relaxation dynamics around quasi-homogeneous configurations.
    They are regime-dependent approximations, valid only under
    weak-field and slow-variation conditions.

  \paragraph{Descriptive hierarchy.}
    The relational formulation is primary; the geometric description
    is a secondary and regime-limited representation.
    Breakdown of geometry therefore reflects loss of applicability
    of the effective language, not inconsistency of the underlying
    relational structure.
